% Backup slides --- shown only in Q&A.

\begin{frame}[plain]
  \centering\Huge Appendices
\end{frame}

\begin{frame}{Magnetar context}
  \centering
  \includegraphics[width=0.8\linewidth]{gertsenshtein_zeldovich_star}
  \vspace{0.4em}

  Magnetar field $B\sim 10^{15}$\,G, scale $\sim 10$\,km: $P_{g\to\gamma}\sim 10^{-10}$.
\end{frame}

\begin{frame}{Full PGT $+$ EM Lagrangian}
  \placeholder{Full Lagrangian with each term labelled in English: Einstein-Cartan, torsion kinetics, curvature-squared, Maxwell, non-minimal couplings.}
\end{frame}

\begin{frame}{Higher-derivative terms: $\Rt^2$ \& $F^4$ via perturbative reduction}
  \scriptsize
  \textbf{Problem.} $\Rt^2$ (and Euler--Heisenberg $(F^2)^2$) produce 4th-order time
  derivatives in the EOMs --- generically Ostrogradsky-ghost modes.  Cannot integrate
  the full theory directly.

  \vspace{0.4em}
  \textbf{Parker--Simon iterative reduction (1993).}  Treat $b_5$ as small;
  the full EOM splits as $E_{\rm base}[y] + b_5\,E_{\rm corr}[y]=0$, where
  $E_{\rm corr}$ contains the 4th-order $\Rt^2$ pieces.  Expand
  $y(k,t) = y^{(0)} + b_5\,y^{(1)} + \cdots$ and collect by order:
  \[
    E_{\rm base}\bigl[y^{(0)}\bigr] = 0,\qquad
    E_{\rm base}\bigl[y^{(n)}\bigr] = -E_{\rm corr}\bigl[y^{(n-1)}\bigr],
    \quad n\ge 1.
  \]
  \emph{The 4th-order operator never acts on the unknown.}  $E_{\rm corr}$ evaluates
  on the previous order's known trajectory, becoming a time-dependent source for the
  next order.  The 4th-order operator is never integrated; ghost eigenvalues never enter the
  propagator.

  The sourced equation is solved analytically by \emph{Duhamel's principle}: convolve
  the base propagator $e^{A(t-\tau)}$ against the source, evaluated analytically in
  the order-$0$ eigenbasis, no numerical quadrature.

  \vspace{0.4em}
  \textbf{Constraint-promotion barrier.}
  At $b_5=0$ (Einstein--Cartan), some metric components are \emph{pure algebraic
  constraints}: non-dynamical, carrying no degrees of freedom.  The $\Rt^2$ correction
  promotes them to propagating fields --- the number of propagating modes
  jumps discontinuously as $b_5\to 0$. No published perturbative Hamiltonian recipe handles
  this case; energy measurements break down across $b_5=0$.
\end{frame}

\begin{frame}{\textsc{tidal} architecture deep-dive}
  \placeholder{CoefficientEvaluator cache levels, operator library, constraint-IC solver, RHS evaluator, state layout.}
\end{frame}

\begin{frame}{Modal solver: two structural fixes}
  \footnotesize
  After Fourier decomposition each spatial mode $k$ satisfies
  $\dot{y}(t,k)=M(k)\,y(t,k)$, solved by $y(t,k)=e^{M(k)t}y_0(k)$.
  Two main problems arise for gauge theories.

  \vspace{0.3em}
  \textbf{(i) Ill-conditioned eigenvectors.}\\
  The standard route $V\,\mathrm{diag}(e^{\lambda t})\,V^{-1}y_0$ requires inverting the
  eigenvector matrix $V$.  Gauge theories have degenerate and near-zero eigenvalues;
  $\mathrm{cond}(V)$ reaches $10^{15}$--$10^{17}$, making $V^{-1}$ numerically catastrophic.\\[0.15em]
  \textbf{Fix:} compute $e^{Mt}y_0$ directly via Pad\'e approximant --- no eigenvectors, machine-precision output.

  \vspace{0.3em}
  \textbf{(ii) Constraint fields (no time derivative).}\\
  ``Constraints'' enter algebraically, not dynamically.
  Per mode, partition into dynamic $d$ and constrained $c$:
  \[
    \begin{pmatrix}\dot{d}\\0\end{pmatrix}
    =
    \begin{pmatrix}S_{dd}&S_{dc}\\S_{cd}&S_{cc}\end{pmatrix}
    \begin{pmatrix}d\\c\end{pmatrix}
    \;\Longrightarrow\;
    c = -S_{cc}^{-1}S_{cd}\,d,\qquad
    \dot{d} = \underbrace{\bigl(S_{dd}-S_{dc}S_{cc}^{-1}S_{cd}\bigr)}_{M_{\rm red}}\,d.
  \]
  Evolve $M_{\rm red}$ by matrix exponential; recover $c$ algebraically at the end.
\end{frame}
